\documentclass[12pt,a4paper]{article}

\usepackage[utf8]{inputenc}
\usepackage[T2A]{fontenc}
\usepackage[russian]{babel}

\usepackage{amsmath,amssymb}
\usepackage{geometry}
\geometry{margin=2.2cm}
\usepackage{microtype}
\usepackage{parskip}

\newcommand{\Z}{\mathbb{Z}}
\DeclareMathOperator{\gcdop}{gcd}
\DeclareMathOperator{\ord}{ord}

\begin{document}
\section*{\centering Лекция 1}
\section{Теорема Эйлера и функция Эйлера}

\textbf{Теорема Эйлера.} Если $a\in(\Z/n)^{\ast}$ (то есть $\gcdop(a,n)=1$), то
\[
a^{\varphi(n)} \equiv 1 \pmod n.
\]

\textbf{Замечания.}
\begin{itemize}
  \item $\Z/n$ — кольцо.
  \item $(\Z/n)^{\ast}$ — группа обратимых элементов.
  \item $\varphi(n)$ — число элементов $x\in\{1,\dots,n-1\}$, взаимно простых с $n$.
\end{itemize}

\textbf{Замечание.} $\varphi(n)=|(\Z/n)^{\ast}|$.

\textbf{Критерий обратимости.} Для $m\in/n$:
\[
m\in(\Z/n)^{\ast} \iff \gcdop(m,n)=1.
\]
Идея: $\gcdop(m,n)=1 \Rightarrow \exists x,y\in: mx+ny=1 \Rightarrow mx\equiv 1\pmod n$.
Обратно: $mc\equiv 1\pmod n \Rightarrow mc-1=kn \Rightarrow \gcdop(m,n)=1$.

\textbf{Примеры.}
\[
\varphi(p)=p-1,\qquad
\varphi(p^{s})=p^{s}-p^{s-1}.
\]
Если $\gcdop(m,k)=1$, то $\varphi(mk)=\varphi(m)\varphi(k)$.

\textbf{Доказательство теоремы Эйлера (идея).}
Пусть $x_1,\dots,x_{\varphi(n)}$ — все элементы $(\Z/n)^{\ast}$.
Тогда $ax_1,\dots,ax_{\varphi(n)}$ — просто перестановка $(\Z/n)^{\ast}$, значит
\[
(ax_1)\cdots(ax_{\varphi(n)}) \equiv x_1\cdots x_{\varphi(n)} \pmod n
\Rightarrow a^{\varphi(n)}\equiv 1\pmod n.
\]

\textbf{Следствие (теорема Ферма).} Если $p$ — простое и $a\in(\Z/p\Z)^{\ast}$, то
\[
a^{p-1}\equiv 1 \pmod p.
\]

\textbf{Порядок элемента.}
Если $a\in(\Z/n)^{\ast}$, то наименьшее $k\in \mathbb{N}$, для которого $a^k\equiv 1\pmod n$,
называется порядком $a$; обозначение $k=\ord a$.

\textbf{Примеры.}
\[
(\Z/2)^{\ast}=\{1\},\quad
(\Z/3)^{\ast}=\{1,2\},\quad
(\Z/4)^{\ast}=\{1,3\},\quad
(\Z/5)^{\ast}=\{1,2,3,4\},\quad
(\Z/6)^{\ast}=\{1,5\}.
\]
\[
(\Z/7)^{\ast}\cong \Z_6,\qquad
(\Z/8)^{\ast}=\{1,3,5,7\}\cong \Z_2\times \Z_2.
\]

\section{Циклические группы, порядок, линейные сравнения}

\textbf{Определение.} Группа $G$ называется циклической, если она порождается одним элементом:
$G=\langle g\rangle$.

\textbf{Лемма про порядок (кратко).}
Пусть $a\in(\Z/n)^{\ast}$ и $k=\ord(a)$, то есть $a^k\equiv 1\pmod n$.
Тогда:
\[
a^m\equiv 1\pmod n \Rightarrow k\mid m,
\qquad\text{в частности}\qquad
\ord(a)\mid \varphi(n).
\]

\textbf{Упражнение.} Сравнение
\[
ax\equiv b \pmod m
\]
имеет решение $\Leftrightarrow \gcdop(a,m)\mid b$.
Если это так, то существует ровно $\gcdop(a,m)$ решений по модулю $m$.

\section{Теорема Вильсона и КТО}

\textbf{Первый тест на простоту (Вильсон).}
\[
p\ \text{простое} \iff (p-1)!\equiv -1\pmod p.
\]
(Идея: в произведении остаются только решения $x^2\equiv 1\pmod p$, то есть $x\equiv\pm1$.)

\textbf{Китайская теорема об остатках.}
Пусть $m_1,\dots,m_r\in \mathbb{N}$ попарно взаимно просты.
Для любых $a_1,\dots,a_r\in$ система
\[
\begin{cases}
x\equiv a_1 \pmod{m_1}\\
\vdots\\
x\equiv a_r \pmod{m_r}
\end{cases}
\]
имеет решение, и оно единственно по модулю $m=\prod_{i=1}^r m_i$.

\textbf{Конструкция.}
Пусть $m=\prod m_i$, $m_i'=\dfrac{m}{m_i}$.
Тогда $\gcdop(m_i',m_i)=1$, поэтому $\exists \alpha_i\in/m_i\Z$:
\[
\alpha_i m_i' \equiv 1\pmod{m_i}.
\]
При этом автоматически $\alpha_i m_i'\equiv 0\pmod{m_j}$ при $j\neq i$.
Решение:
\[
x_0=\sum_{i=1}^r a_i\,\alpha_i\,m_i'.
\]
Если $x_0$ и $\tilde x_0$ — два решения, то $m_i\mid(x_0-\tilde x_0)$ для всех $i$,
а значит $m\mid(x_0-\tilde x_0)$ и решения совпадают по модулю $m$.

\textbf{Культурная переформулировка.}
Отображение
\[
\Phi:\Z/m\Z \to \Z/m_1\Z\times\cdots\times \Z/m_r\Z,\qquad
\Phi(\bar a)=(\bar a\bmod m_1,\dots,\bar a\bmod m_r)
\]
— изоморфизм колец (гомоморфизм по сложению и умножению).

\section{Следствие КТО для обратимых и формула для $\varphi(n)$}

\textbf{Отсюда.}
\[
(\Z/m\Z)^{\ast}\cong (\Z/m_1\Z)^{\ast}\times\cdots\times(\Z/m_r\Z)^{\ast}.
\]
В частности, при $\gcdop(k,s)=1$:
\[
\varphi(ks)=\varphi(k)\varphi(s).
\]

\textbf{Как вычислять $\varphi(n)$.}
Если $n=\prod_{i} p_i^{\alpha_i}$, то
\[
\varphi(n)=\prod_i \varphi(p_i^{\alpha_i})
=\prod_i \left(p_i^{\alpha_i}-p_i^{\alpha_i-1}\right).
\]

\section{Чевалле--Уорнинг}

Пусть $K=\mathbb{F}_q$ — конечное поле из $q=p^s$ элементов.
Факт: $K^{\ast}$ — циклическая группа порядка $(q-1)$.

\textbf{Лемма.} Для $u\in \mathbb{N}$ положим
\[
S(u)=\sum_{x\in K} x^u.
\]
Тогда
\[
S(u)=
\begin{cases}
-1, & (q-1)\mid u,\\
0, & \text{иначе}.
\end{cases}
\]
Замечания: $u=0\Rightarrow S(0)=q\equiv 0\pmod p$; если $(q-1)\nmid u$, то
берут $z\in K^{\ast}$ с $z^u\neq 1$ и получают $(1-z^u)S(u)=0$.

\textbf{Теорема Чевалле--Уорнинга.} Пусть $f_1,\dots,f_m\in K[x_1,\dots,x_n]$ и
\[
\sum_{i=1}^m \deg(f_i) < n.
\]
Рассмотрим множество решений
\[
V=\{x\in K^n \mid f_1(x)=\cdots=f_m(x)=0\}.
\]
Тогда
\[
|V|\equiv 0\pmod p.
\]
(аффинное алгебраическое многообразие.)

\textbf{Идея доказательства.}
Рассматривают
\[
P(x)=\prod_{i=1}^m \bigl(1-f_i(x)^{\,q-1}\bigr).
\]
Если $x\in V$, то $P(x)=1$.
Если $x\notin V$, то найдётся $i$ с $f_i(x)\neq 0$, тогда $f_i(x)^{q-1}=1$ и $P(x)=0$.
Значит
\[
\sum_{x\in K^n} P(x)=|V|.
\]
Далее раскладывают $P$ в сумму мономов и применяют лемму про суммы степеней по каждой переменной,
используя ограничение $\deg(P)<n(q-1)$ (из $\sum \deg f_i < n$).

\section{Интермеццо: теорема Гильберта о базисе}

Пусть $K$ — поле, $K[x_1,\dots,x_n]$.
Тогда любой идеал $M\subset K[x_1,\dots,x_n]$ конечно порождён:
\[
\exists f_1,\dots,f_N\in K[x_1,\dots,x_n] \ \forall f\in M\ \exists h_1,\dots,h_N\in K[x_1,\dots,x_n]:
\quad
f=\sum_{i=1}^N h_i f_i.
\]
(кольцо многочленов над полем Нётерово.)

\section{RSA}

\textbf{Дано.}
Берём «случайные» большие простые $p,q$ и $n=pq$.
Выбираем $e$ так, что
\[
e<\varphi(n)=(p-1)(q-1),\qquad \gcdop(e,\varphi(n))=1.
\]
Находим $d$:
\[
de\equiv 1 \pmod{\varphi(n)}.
\]

\textbf{Шифрование и расшифрование.}
\[
m\in \Z/n,\qquad E(m)=m^e\pmod n,\qquad D(c)=c^d\pmod n.
\]

\textbf{Ключи.}
Открытый ключ $(e,n)$ знают все. Секретный ключ $(d,n)$ знает владелец.

\textbf{Корректность.}
Нужно: $E(D(m))\equiv m\pmod n$.
Достаточно показать по модулю $p$ и по модулю $q$:
\[
E(D(m))\equiv m\pmod p,\qquad E(D(m))\equiv m\pmod q.
\]
Если $m\equiv 0\pmod p$, то очевидно.
Если $m\not\equiv 0\pmod p$, то из $de\equiv 1\pmod{(p-1)(q-1)}$ следует
$de=1+c(p-1)(q-1)$, и тогда
\[
m^{de}=m^{1+c(p-1)(q-1)} \equiv m\cdot (m^{p-1})^{c(q-1)}\equiv m \pmod p
\]
по малой теореме Ферма. Аналогично по модулю $q$.

\section{Упражнение: $x^2\equiv -1\pmod p$ и пример обратного элемента}

\textbf{Упражнение.}
\[
x^2\equiv -1\pmod p \iff p\equiv 1\pmod 4.
\]
Связка через МТФ/Вильсона:
\[
x^{p-1}\equiv 1\pmod p,\qquad
x^{p-1}=(x^2)^{\frac{p-1}{2}}\equiv (-1)^{\frac{p-1}{2}},
\]
и анализ $(p-1)!\equiv -1\pmod p$ через попарные множители.

\textbf{Пример: найти $5^{-1}\pmod{36}$.}
Решаем $5x\equiv 1\pmod{36}$, то есть диофантово
\[
5x+36y=1.
\]
Из цепочки:
\[
36=5\cdot 7+1 \Rightarrow 1=36-5\cdot 7
\Rightarrow x\equiv -7\equiv 29\pmod{36}.
\]

Альтернатива через Эйлера:
\[
a^{-1}\equiv a^{\varphi(n)-1}\pmod n,\qquad \varphi(36)=12,\qquad 5^{11}\equiv 29\pmod{36}.
\]
Также:
\[
(\Z/36\Z)^{\ast}=\{1,5,7,11,13,17,19,23,25,29,31,35\}.
\]

\section{Тесты на простоту, тест Люка, числа Кармайкла}

\textbf{Тесты на простоту}
\begin{enumerate}
  \item $p$ простое $\Rightarrow (p-1)!\equiv -1\pmod p$.
  \item $p$ простое $\Rightarrow$ все решения $x^2\equiv 1\pmod p$ это $x\equiv \pm 1$.
  \item $p$ простое $\Rightarrow \forall a,\ 1\le a\le p-1:\ a^{p-1}\equiv 1\pmod p$ (МТФ).
\end{enumerate}
Верно и обратное:
если $\forall a\in\{1,\dots,n-1\}$ выполнено $a^{n-1}\equiv 1\pmod n$, то $n$ — простое.

\textbf{Теорема Люка.}
Если найдётся $a\in(\Z/n)^{\ast}$ такое, что
\[
a^{n-1}\equiv 1\pmod n
\quad\text{и}\quad
a^{\frac{n-1}{p}}\not\equiv 1\pmod n\ \ \forall\ p\ \text{простое, }p\mid(n-1),
\]
то $n$ — простое.

\textbf{Идея доказательства.}
Если $n$ не простое, то $\varphi(n)<n-1$.
Для $a\in(\Z/n)^{\ast}$ верно $a^{\varphi(n)}\equiv 1$.
Пусть $d=\gcdop(n-1,\varphi(n))$, тогда $a^d\equiv 1$, и это приводит к противоречию
с условием $a^{\frac{n-1}{p}}\not\equiv 1$.

\textbf{Числа Кармайкла.}
Определение: $m\in \mathbb{N}$ — число Кармайкла, если $m$ составное и
\[
a^m \equiv a \pmod m \quad \forall a\in.
\]
Замечание: минимальное $m=561$; таких чисел бесконечно много (доказано в 1992).

\textbf{Утверждение (критерий Корсельта).}
$m$ — число Кармайкла $\iff$
\begin{itemize}
  \item $m$ свободно от квадратов;
  \item для каждого простого делителя $p\mid m$ выполнено $(p-1)\mid(m-1)$.
\end{itemize}

\textbf{Мораль.}
Числа Кармайкла показывают «слабость» теста Ферма.\\
\textbf{Вывод:} тест Миллера--Рабина.

\end{document}